\clearpage\thispagestyle{empty}
\pdfbookmark[0]{Sources for the Heliocentric Rules (1802)}{vorschriften1802-sources}
\begin{center}\textbf{Sources for the Heliocentric Rules (1802)}\end{center}
\noindent This AI translation follows Carl Friedrich Gauss’s German text in \textit{Werke}, vol.~VI (Göttingen, 1874), pp.~87--93. The first printing appeared in the June 1802 issue of \textit{Monatliche Correspondenz zur Beförderung der Erd- und Himmels-Kunde}, edited by Franz Xaver von Zach, vol.~5, pp.~540--546, under the name Dr.~Gauss in Braunschweig. \href{https://gdz.sub.uni-goettingen.de/id/PPN236007467}{Collected German text}. \href{https://www.digitale-sammlungen.de/en/view/bsb10538597?page=556}{First printing}.
\par\smallskip\noindent The journal denotes the Sun’s longitude by \(V\), whereas the collected text uses \(\odot\). In worked calculation \(6^{0}\), the journal prints an \(n\) after the complement of \(\log\cos i\) (p.~544). The collected text gives \(0.0075025\) without that mark (p.~91). The edition here follows the collected reading.
\par\smallskip\noindent A reprint is in Appendix IV, pp.~42--45, of \href{https://gdz.sub.uni-goettingen.de/id/PPN60187157X}{Carl Haase’s German translation of \textit{Theoria motus}} (Hannover, 1865). For Gauss’s related treatment and reference back to the journal article, see Article~74, pp.~89--91, in \href{https://www2.karlin.mff.cuni.cz/~tuma/Aplikace21/theoryofmotionof00gausuoft.pdf#page=113}{Charles Henry Davis’s English translation of \textit{Theoria motus}} (Boston, 1857).
\par\smallskip\noindent Translation and typesetting produced by AI.
\clearpage
